\chapter{驱动、执行器与传感}
\label{ch:drive-systems}

% ════════════════════════════════════════════════════════════════
%  机械臂建模、规划与控制部 · 驱动/执行器/传感章。
%  集大成源：Siciliano Ch5（执行器与传感，主线）+ Craig §9.9（单关节模型/
%  有效惯量/结构共振带宽）+ Spong §6.1.1（反映惯量/弹性关节）+ Paul §6.3（执
%  行器惯量降低结构依赖的实证）。
%  承上：ch:arm-dyn 给「负载是什么」（M,C,g）。
%  本章：给「谁来驱动负载、驱动器是什么模型」。
%  启下：ch:arm-ctrl 的关节控制需要执行器模型；P4 sec:ctrl-robot。
%  记号纪律：电气量用工程惯例（V_a,I_a,R_a,L_a,k_e,k_t），机械侧与全书一致
%  （J_m 电机惯量、k_r 减速比、tau 力矩）。
% ════════════════════════════════════════════════════════════════

\begin{practice}[前置自测]
开始之前，先试答下列问题。若已能流畅作答，可只速读本章的速查卡与符号表；若卡住，正是本章要为你解决的。
\begin{enumerate}
  \item 上一章（\cref{ch:arm-dyn}）把“给定运动求所需关节力矩 $\boldsymbol{\tau}$”讲透了。可是 $\boldsymbol{\tau}$ 不会凭空出现——它由电机经减速器产生。一台直流电机，你能写出“施加电枢电压 $\to$ 输出轴力矩”这条因果链的微分方程吗？反电动势在其中扮演什么角色？
  \item 同一台电机，配上不同的驱动电路，可以表现为“给电压、转速正比于电压”的\emph{速度源}，也可以表现为“给指令、力矩正比于指令、与转速几乎无关”的\emph{力矩源}。这两种行为的电气区别在哪里？计算力矩控制（\cref{sec:ctrl-robot-ctc}）需要哪一种？
  \item 工业机械臂的减速比 $k_r$ 动辄几十到几百。一个反复被引用的结论是：减速比把负载惯量“反映”到电机侧时要除以 $k_r^2$。请先猜：为什么是平方、而不是一次方？这个平方因子如何解释“高减速臂的各关节近似解耦、可当独立单关节来控”？
  \item 把电机、减速器、负载串起来，没有任何元件是绝对刚性的——轴会扭、齿会弹、谐波减速器的柔轮本就靠弹性变形工作。这条柔性链会冒出一个\emph{结构共振峰}。为什么控制器增益开大、闭环带宽推高，会“撞上”这个共振峰而失稳？工程上 $\omega_n\le\tfrac12\omega_{\mathrm{res}}$ 这条经验律是怎么来的？
  \item 编码器测的是\emph{电机轴}角度还是\emph{关节}角度？在有减速器、有柔性的真实关节上，这二者相差一个 $k_r$ 倍加上一个扭转变形——这个差别会给“电机侧位置环”埋下什么隐患？
\end{enumerate}
\end{practice}

\section*{本章目标}
学完本章，你应当能够：
\begin{enumerate}
  \item \textbf{建立}永磁直流电机的电气--机械耦合模型：电枢方程 $V_a=(R_a+sL_a)I_a+k_e\Omega_m$、转矩方程 $C_m=k_tI_a$，并由此画出方块图、写出“控制电压 $\to$ 转角”的二阶传递函数（\cref{sec:ds-dcmotor}）；
  \item \textbf{区分}速度型与力矩型驱动器架构：电流环增益如何决定驱动器是\emph{速度控制发电机}（$k_i=0$，强扰动抑制）还是\emph{力矩控制发电机}（$Kk_i\gg R_a$，力矩与转速解耦），并知道独立关节控制与集中式控制各偏好哪一种（\cref{sec:ds-drives}）；
  \item \textbf{推导}减速比的“反映律”：负载惯量/黏性摩擦按 $1/k_r^2$、反作用（负载）力矩按 $1/k_r$ 折算到电机侧，得有效惯量 $J_{\mathrm{eff}}=J_m+J_\ell/k_r^2$（电机侧）或 $k_r^2 J_m+J_\ell$（负载侧），并据此说清高减速比为何“隐藏”了关节间的非线性耦合（\cref{sec:ds-gear}）；
  \item \textbf{分析}谐波减速器/行星齿轮的传动特性，建立弹性关节的二自由度模型 $J_\ell\ddot{q}_\ell+k(q_\ell-q_m)+\cdots=0,\ J_m\ddot{q}_m+k(q_m-q_\ell)=u$，并理解背隙与柔性对控制的代价（\cref{sec:ds-flex}）；
  \item \textbf{说清}液压执行器的输入/输出模型与“为何重载臂偏爱液压”，以及它与电气驱动器在方块图上的形式类比（\cref{sec:ds-hydraulic}）；
  \item \textbf{选用}本体感知与外感知传感器：增量/绝对编码器、旋转变压器、测速机、应变片式力/力矩传感器、腕力传感器的标定矩阵（\cref{sec:ds-sensors}）；
  \item \textbf{用第一性原理判断控制带宽上限}：把电机$+$减速器$+$负载视为带柔性的三阶链，估算结构共振 $\omega_{\mathrm{res}}=\sqrt{k/m_{\mathrm{eff}}}$，并理解经验安全裕度 $\omega_n\le\tfrac12\omega_{\mathrm{res}}$ 的来由与应用（\cref{sec:ds-bandwidth}）。
\end{enumerate}

\subsection*{本章知识导航}
本章在\emph{动力学}与\emph{控制}之间架桥。\cref{ch:arm-dyn} 回答“负载是什么”——它把机械臂折成方程 $\mathbf{M}(\mathbf{q})\ddot{\mathbf{q}}+\mathbf{C}\dot{\mathbf{q}}+\mathbf{g}=\boldsymbol{\tau}$，左边是“要驱动的负载”，右边的 $\boldsymbol{\tau}$ 却被默认“想给多少就有多少”。本章戳破这个理想化：$\boldsymbol{\tau}$ 来自\emph{电机经减速器}，电机有自己的电气惯性与饱和，减速器有自己的反映效应与柔性。理解了“谁来驱动负载、用什么模型驱动”，才能在 \cref{ch:arm-ctrl} 写出真正可落地的关节控制器——尤其是那条“控制器带宽不能高过结构共振一半”的硬约束，它直接决定了控制增益的天花板。

\begin{center}
\begin{forest} navtreeLR
[{驱动 / 执行器 / 传感}
  [{直流电机模型\\电枢$+$反电动势$+$转矩}
    [{二阶传递函数 $M(s)$}] ]
  [{驱动器架构\\速度型 vs 力矩型}
    [{电流环增益决定}] ]
  [{减速比反映\\$J_{\mathrm{eff}}=J_m+J_\ell/k_r^2$}
    [{隐藏耦合 / 线性化}] ]
  [{谐波/行星·柔性\\弹性关节 2-DOF}
    [{背隙 / 共振}] ]
  [{液压执行器\\重载 / 高力矩}]
  [{传感\\编码器/力矩/IMU}]
  [{带宽与共振\\$\omega_n\le\tfrac12\omega_{\mathrm{res}}$}
    [{增益天花板}] ]
]
\end{forest}
\end{center}

\subsection*{前置桥接}
本章复用本书已建的地基，遇到时一律交叉引用回去、不重证：
\begin{itemize}
  \item \textbf{机械臂动力学}：负载方程 $\mathbf{M}(\mathbf{q})\ddot{\mathbf{q}}+\mathbf{C}\dot{\mathbf{q}}+\mathbf{g}=\boldsymbol{\tau}$ 与构型相关惯量 $\mathbf{M}(\mathbf{q})$ 来自 \cref{ch:arm-dyn}（P4 \cref{sec:ctrl-robot-dyn} 亦给出结论）；本章把“单个关节看见的有效惯量”接到这条主线上。
  \item \textbf{运动学/雅可比}：编码器测的电机角经减速比与逆运动学才换成关节角与末端位姿，几何链来自 \cref{ch:arm-kin}、\cref{sec:ak-jacobian}。
  \item \textbf{控制导论}：本章末端给出的单关节二阶/三阶模型，正是 P4 \cref{sec:ctrl-robot} 独立关节控制（PD$+$前馈）与 LQR（\cref{sec:ctrl-lqr}）的\emph{被控对象}；那里讲“怎么控”，本章讲“控的是什么、增益上限在哪”。
\end{itemize}

% ════════════════════════════════════════════════════════════════
\section{关节驱动系统总览}
\label{sec:ds-overview}

机械臂每个关节的运动，由一套\textbf{驱动系统}（actuating system）实现。沿功率流，它一般由四级串联组成：

\begin{equation}
\underbrace{\text{电源}}_{P_p}\;\longrightarrow\;\underbrace{\text{功率放大器}}_{P_a}\;\longrightarrow\;\underbrace{\text{伺服电机}}_{P_m}\;\longrightarrow\;\underbrace{\text{传动（减速器）}}_{P_u}\,.
\label{eq:ds-powerflow}
\end{equation}

每一级把上一级的功率\emph{转换}并向下传递，转换中各级分别耗散掉 $P_{da},P_{ds},P_{dt}$（放大、电机、传动的损耗）。功率在任何一级都可写成“一个流量 $\times$ 一个力量”：电气功率 $=$ 电压 $\times$ 电流，机械功率 $=$ 力矩 $\times$ 角速度，液压功率 $=$ 压力 $\times$ 流量。控制律产生的（通常是电气的）信号功率 $P_c$ 经放大器调制后注入这条链；末端 $P_u$ 就是关节实际拿到的机械功率。

\begin{insight}[选型从“末端要多大机械功率”倒推]
\label{ins:ds-power-backward}
设计一套关节驱动，应当\emph{从负载侧倒着选}：先由任务要求的关节力矩与角速度定出 $P_u$，再上溯决定电机、放大器、电源的规格。这与控制器设计“正着走”的视角互补——控制器假定执行器能提供所需 $\boldsymbol{\tau}$，而执行器选型恰恰是去保证“这个假定成立”。\cref{ch:arm-traj} 把力矩约束塞进轨迹时间参数化、\cref{ch:arm-prac} 用真力矩约束跑 TOPP，背后正是这条“$P_u$ 倒推”链。
\end{insight}

\paragraph{为何需要传动。}\;关节运动的典型需求是\emph{低速、大力矩}；而伺服电机的最优工作区恰恰相反——\emph{高速、小力矩}。于是必须插入一级\textbf{传动}（gear）把电机的高速小力矩“变速”成关节的低速大力矩。除变速外，传动还能改变运动\emph{性质}（电机的旋转经丝杠变为移动关节的平动），并能把电机\emph{上移}：若把若干电机装到机座，整臂总重下降、功率重量比上升。常见传动有正齿轮（变轴/平移作用点）、丝杠（旋转$\to$平动，驱动移动关节，常用滚珠丝杠预紧以增刚减背隙）、同步带与链（把电机远置）。代价是引入\emph{弹性与背隙}——这正是 \cref{sec:ds-flex} 与 \cref{sec:ds-bandwidth} 的主题。

\paragraph{直驱（direct drive）。}\;若电机力矩特性允许，可不经传动直连关节，从而\emph{根除}传动弹性与背隙。代价是：少了减速比的“反映”，动力学里的非线性耦合项不再可忽略（见 \cref{ins:ds-hide-coupling}），需要更复杂的控制（如计算力矩，\cref{sec:ctrl-robot-ctc}）；同等力矩的直驱电机又大又贵。故直驱在工业臂上尚不普及，多见于对带宽要求极高的科研臂。

\paragraph{电机类型。}\;按输入功率 $P_a$ 的性质，电机分气动、液压、电动三类。气动难精确控制（流体可压缩误差不可避），多只用于夹爪开合一类两位运动。机器人最常用\textbf{电气伺服电机}，其中\emph{永磁直流（DC）}与\emph{无刷直流（BLDC）}因控制灵活性最受欢迎。二者在\emph{建模}上可用同一组微分方程描述（无刷电机配换相模块与位置传感器后，电气行为等价于电刷处于与励磁磁通成 $\pi/2$ 的永磁直流机），故下文统一用“直流电机模型”。无刷的优势在去掉机械换相（电刷$+$换向器）带来的电气/机械损耗与磨损，并把绕组移到定子利于散热、转子更紧凑惯量更低。重载场合则用\emph{液压伺服电机}（\cref{sec:ds-hydraulic}）。

% ════════════════════════════════════════════════════════════════
\section{直流电机的电气--机械模型}
\label{sec:ds-dcmotor}

这一节把“施加电枢电压 $\to$ 输出轴转角”的因果链建成方块图与传递函数。它是后面一切的地基：驱动器架构（\cref{sec:ds-drives}）是在这张方块图上加电流反馈环，单关节控制（\cref{sec:ds-bandwidth}）是再串上负载与柔性。

\subsection{两条平衡方程}
\label{sec:ds-balances}

直流电机的物理本质是\emph{洛伦兹力}：通电绕组（电流 $i_a$）在定子磁场 $\mathbf{B}$ 中受力，合成出与电流成正比的电磁\textbf{转矩}\cite{craig2005introduction}：

\begin{equation}
C_m = k_t\,I_a,
\label{eq:ds-torque}
\end{equation}
其中 $k_t$ 是\textbf{转矩常数}。反过来，电机转动时绕组切割磁力线，像发电机一样在电枢两端产生一个与转速成正比的\textbf{反电动势}（back-EMF）：

\begin{equation}
V_g = k_e\,\Omega_m,
\label{eq:ds-backemf}
\end{equation}
$k_e$ 是\textbf{反电动势常数}（电压常数，文献亦记 $k_v$）。在 SI 单位制下、对补偿良好的电机，二者数值相等：$k_t=k_e$（量纲分别为 $\mathrm{N\cdot m/A}$ 与 $\mathrm{V\cdot s/rad}$，数值相同源于能量守恒——电功率 $V_g I_a$ 恰等于机械功率 $C_m\Omega_m$）。

\textbf{电气平衡}（电枢回路）：电枢由电压源 $V_a$、电阻 $R_a$、电感 $L_a$ 与反电动势 $V_g$ 串联构成，一阶 RL 电路。在复频域（拉氏变量 $s$）：

\begin{equation}
V_a = (R_a + sL_a)\,I_a + V_g = (R_a + sL_a)\,I_a + k_e\,\Omega_m.
\label{eq:ds-electrical}
\end{equation}

\textbf{机械平衡}（电机轴）：驱动转矩 $C_m$ 克服转子惯量 $I_m$、黏性摩擦 $F_m$，并平衡折算到电机轴的负载反作用力矩 $C_l$：

\begin{equation}
C_m = (sI_m + F_m)\,\Omega_m + C_l.
\label{eq:ds-mechanical}
\end{equation}

\begin{derivation}[从两条平衡到统一传递函数]
把 \cref{eq:ds-torque,eq:ds-backemf,eq:ds-electrical,eq:ds-mechanical} 串起来。消去 $I_a$ 与 $C_m$：由 \cref{eq:ds-electrical} 得 $I_a=\frac{V_a-k_e\Omega_m}{R_a+sL_a}$，代入 $C_m=k_tI_a$ 再代入 \cref{eq:ds-mechanical}，得电机轴角速度对电压、负载力矩的响应。\emph{降阶}处理：电气时间常数 $L_a/R_a$ 远小于机械时间常数 $I_m/F_m$（前者 $10^{-3}\,$s 量级，后者 $10^{-1}\,$s 量级），故令 $L_a\approx0$。又因机械黏性摩擦远小于“电气摩擦” $F_m\ll k_ek_t/R_a$（这是直流电机的一个关键不等式：反电动势提供的等效阻尼远大于轴承黏滞），化简后“控制输入 $\to$ 电机转角 $\theta_m$”可\emph{统一}写成
\begin{equation}
M(s)=\frac{\Theta_m(s)}{V_c(s)}=\frac{k_m}{s\,(1+sT_m)},
\label{eq:ds-Ms}
\end{equation}
其中 $k_m$ 为等效增益、$T_m$ 为机电时间常数（其具体表达式因驱动器是速度型还是力矩型而异，见 \cref{sec:ds-drives}）。注意分母里的因子 $s$：电机本质是个\emph{积分器}（“转速积分得转角”），这是关节位置环必须围绕的核心。
\end{derivation}

\begin{note}[功率放大器：一个可忽略的快极点]
放大器把控制电压 $V_c$ 放大成电枢电压 $V_a$，其传递函数 $\frac{V_a}{V_c}=\frac{G_v}{1+sT_v}$。采用 $10$--$100\,$kHz 的 PWM（脉宽调制）开关，放大器时间常数 $T_v\sim10^{-5}$--$10^{-4}\,$s，远小于系统其他时间常数，分析时常令 $T_v\approx0$，只保留增益 $G_v$。PWM 晶体管放大器功率转换效率 $>0.9$、功率增益 $\sim10^6$；驱动永磁直流机的是 DC--DC 斩波器，驱动无刷的是 DC--AC 逆变器。
\end{note}

\begin{pitfall}[$k_t=k_e$ 数值相等不等于量纲相同]
初学者常把“转矩常数等于反电动势常数”误读成“它们是同一个量”。它们\emph{量纲不同}（$\mathrm{N\cdot m/A}$ vs $\mathrm{V\cdot s/rad}$），只是在 SI 制下\emph{数值}巧合相等，根因是能量守恒（无损电机里电功率全转机械功率）。换到非 SI 单位、或电机有铁损/磁滞，这个数值相等就破缺，必须分别标定。手册给的 $K_T$ 与 $K_E$ 若数值不等，多半是单位制或损耗所致，别当矛盾。
\end{pitfall}

% ════════════════════════════════════════════════════════════════
\section{驱动器架构：速度型 vs 力矩型}
\label{sec:ds-drives}

同一台电机配不同的电流反馈，可表现为两种迥异的“理想源”。区别全在\textbf{电流环}：在 \cref{eq:ds-Ms} 的方块图上，于放大器与电枢之间插一个电流传感器 $k_i$ 与电流调节器 $C_i(s)$，构成电枢电流的局部反馈环。这一环的回路增益大小，决定了驱动器的对外行为。

\subsection{速度控制发电机（$k_i=0$）}

\emph{不加}电流反馈（$k_i=0$）。利用 $F_m\ll k_ek_t/R_a$，并设 $C_l=0$，由模型直接得

\begin{equation}
\omega_m \approx \frac{G_v}{k_e}\,v_c',
\label{eq:ds-velgen}
\end{equation}
即\emph{角速度正比于控制电压}——驱动器是一台\textbf{速度控制发电机}。把负载力矩 $C_l$ 作为扰动看其影响：

\begin{equation}
\Omega_m=\frac{\frac{1}{k_e}}{1+s\frac{R_aI_m}{k_ek_t}}\,G_vV_c'\;-\;\frac{\frac{R_a}{k_ek_t}}{1+s\frac{R_aI_m}{k_ek_t}}\,C_l.
\label{eq:ds-velgen-tf}
\end{equation}
扰动项系数为 $R_a/(k_ek_t)$，很小——\textbf{扰动抑制好}。此时 \cref{eq:ds-Ms} 中 $k_m=1/k_e,\ T_m=R_aI_m/(k_ek_t)$。

\subsection{力矩控制发电机（$Kk_i\gg R_a$）}

\emph{加}大增益电流反馈（$Kk_i\gg R_a$）。稳态下

\begin{equation}
c_m \approx \frac{k_t}{k_i}\Big(v_c'-\frac{k_e}{G_v}\omega_m\Big),
\label{eq:ds-torquegen}
\end{equation}
因 $G_v$ 很大，第二项可忽略，\emph{驱动力矩与转速几乎无关、正比于指令}——驱动器是一台\textbf{力矩控制发电机}。其转速对扰动的响应

\begin{equation}
\Omega_m=\frac{\frac{k_t}{k_iF_m}}{1+s\frac{I_m}{F_m}}\,V_c'\;-\;\frac{\frac{1}{F_m}}{1+s\frac{I_m}{F_m}}\,C_l,
\label{eq:ds-torquegen-tf}
\end{equation}
扰动项系数为 $1/F_m$，比速度型的 $R_a/(k_ek_t)$ \emph{大得多}——\textbf{扰动抑制差}。此时 $k_m=k_t/(k_iF_m),\ T_m=I_m/F_m$。注意放大器增益 $G_v$ 在力矩型里不再单独出现（它被关进局部环内），只以 $1/k_i$ 因子藏在 $k_m$ 中。

\begin{insight}[选哪一种，取决于上层控制结构]
\label{ins:ds-vel-vs-torque}
\textbf{独立关节控制}（把每个关节当独立二阶系统、关节间耦合视为扰动，见 P4 \cref{sec:ctrl-robot}）希望驱动器\emph{自己}就把扰动力矩压住，故选\textbf{速度型}（$k_i=0$）——\cref{eq:ds-velgen-tf} 的小扰动系数正合需要；保护则靠在反馈路径上加“死区非线性”的电流限幅，而非饱和控制信号。反之，\textbf{集中式（基于模型）控制}（如计算力矩 \cref{sec:ctrl-robot-ctc}、操作空间控制）要亲自计算并下发\emph{力矩}，故要求驱动器是\textbf{力矩型}——它把 $\boldsymbol{\tau}^{\rm cmd}$ 忠实变成轴力矩，控制器自己负责扰动抵消。一句话：\emph{扰动抑制的责任，要么交给驱动器（速度型），要么交给控制律（力矩型），二选一}。现代力控/柔顺/RL 策略几乎都跑在力矩型驱动器上。
\end{insight}

\begin{pitfall}[力矩型不是“真力矩”，是“电流的代理”]
力矩型驱动器的力矩伺服\emph{间接}经电流测量实现（$C_m=k_tI_a$），它伺服的是电流、再乘 $1/k_t$ 当力矩。这意味着：凡是改变“电流$\to$力矩”关系的因素（温升导致磁通漂移、$k_t$ 随转角脉动、换相纹波），都直接污染力矩精度。要真正“力矩免疫于参数漂移”，得用\emph{直接力矩测量}（轴力矩传感器，\cref{sec:ds-sensors}）闭环，而非电流推算。此外，电流环本身的稳定性必须单独保证，尤其用死区电流限幅时。
\end{pitfall}

% ════════════════════════════════════════════════════════════════
\section{减速比的“反映”：有效惯量与隐藏的耦合}
\label{sec:ds-gear}

这是本章最具操作意义的一节。它解释了一个反复出现却常被当作“魔法系数”的事实：减速比 $k_r$ 把负载惯量反映到电机侧要\emph{除以 $k_r^2$}。理解了它，才懂为什么高减速臂能近似当作一串独立的单关节来控。

\subsection{一对正齿轮的力矩平衡}
\label{sec:ds-gearmesh}

考虑电机轴与关节轴经一对正齿轮（半径 $r_m,r$）耦合，理想无背隙。定义\textbf{减速比}

\begin{equation}
k_r=\frac{r}{r_m}=\frac{\vartheta_m}{\vartheta}=\frac{\omega_m}{\omega}\;(\gg 1,\ \text{典型几十到几百}),
\label{eq:ds-ratio}
\end{equation}
即电机转得快、关节转得慢，二者角位移之比恒为 $k_r$（无滑动时 $r_m\vartheta_m=r\vartheta$）。齿间接触力 $f$ 在电机侧产生反作用力矩 $f r_m$、在负载侧产生驱动力矩 $f r$。两侧力矩平衡（$I_m,F_m$ 为电机转子惯量与黏性摩擦，$I,F$ 为负载侧的；齿轮自身惯量已并入相应侧）：

\begin{align}
c_m &= I_m\dot\omega_m + F_m\omega_m + f r_m, \label{eq:ds-mbal-motor}\\
f r &= I\dot\omega + F\omega + c_l. \label{eq:ds-mbal-load}
\end{align}

\begin{derivation}[消去齿间力，折算到电机轴]
由 \cref{eq:ds-mbal-load} 解出 $f=\frac{1}{r}(I\dot\omega+F\omega+c_l)$，代入 \cref{eq:ds-mbal-motor}，并用 \cref{eq:ds-ratio} 把负载侧量换成电机侧量（$\omega=\omega_m/k_r,\ \dot\omega=\dot\omega_m/k_r$，且 $r_m/r=1/k_r$）：
\[
c_m = I_m\dot\omega_m + F_m\omega_m + \frac{r_m}{r}\Big(I\frac{\dot\omega_m}{k_r}+F\frac{\omega_m}{k_r}+c_l\Big)
     = \Big(I_m+\frac{I}{k_r^2}\Big)\dot\omega_m + \Big(F_m+\frac{F}{k_r^2}\Big)\omega_m + \frac{c_l}{k_r}.
\]
\emph{平方因子的来历看得很清}：负载惯量先经一次 $1/k_r$（加速度换算 $\dot\omega=\dot\omega_m/k_r$），再经一次 $1/k_r$（力矩臂比 $r_m/r=1/k_r$），两次相乘即 $1/k_r^2$。而反作用力矩 $c_l$ 只经一次力矩臂比，故只除 $1/k_r$。
\end{derivation}

于是得电机侧的有效参数：

\begin{equation}
\boxed{\;c_m=I_{\rm eq}\dot\omega_m+F_{\rm eq}\omega_m+\frac{c_l}{k_r},\qquad
I_{\rm eq}=I_m+\frac{I}{k_r^2},\quad F_{\rm eq}=F_m+\frac{F}{k_r^2}.\;}
\label{eq:ds-Ieq-motor}
\end{equation}

对偶地，若以\emph{负载侧}变量写（这是控制器更常用的视角，因为期望轨迹在关节空间），把电机力矩 $\tau=k_r c_m$ 折到负载侧、负载角 $\vartheta$ 为坐标：

\begin{equation}
\boxed{\;\tau=(I+k_r^2 I_m)\ddot\vartheta+(F+k_r^2 F_m)\dot\vartheta,\qquad J_{\rm eff}=I+k_r^2 I_m.\;}
\label{eq:ds-Ieq-load}
\end{equation}
注意两式的对偶：电机侧负载惯量\emph{除以} $k_r^2$，负载侧电机惯量\emph{乘以} $k_r^2$——同一个物理事实的两种记账。\cref{eq:ds-Ieq-load} 的 $J_{\rm eff}=I+k_r^2I_m$ 称为“在减速器输出（关节）侧\emph{看见}的有效惯量”，是 \cref{ch:arm-ctrl} 单关节控制的被控惯量\cite{craig2005introduction}。

\begin{example}[Spong 单连杆：用 $I=r^2J_m+J_\ell$ 写运动方程]
\label{ex:ds-spong-single}
取单连杆经齿轮驱动，连杆角 $\theta_\ell$、电机角 $\theta_m=k_r\theta_\ell$。系统总动能 $K=\tfrac12 J_m\dot\theta_m^2+\tfrac12 J_\ell\dot\theta_\ell^2=\tfrac12(k_r^2 J_m+J_\ell)\dot\theta_\ell^2$，即有效惯量 $I=k_r^2 J_m+J_\ell$（与 \cref{eq:ds-Ieq-load} 一致）。势能 $P=Mg\ell(1-\cos\theta_\ell)$。代入欧拉--拉格朗日方程，并把电机力矩 $u=k_r\tau_m$、黏性 $B=k_r B_m+B_\ell$ 折到连杆侧，得
\begin{equation}
I\ddot\theta_\ell+B\dot\theta_\ell+Mg\ell\sin\theta_\ell=u.
\label{eq:ds-spong-eom}
\end{equation}
这恰是把 \cref{ch:arm-dyn} 的拉格朗日法用到“单杆$+$执行器”上的最小例子，也是 P4 独立关节控制的标准被控对象\cite{spong2006robot}。
\end{example}

\subsection{为何高减速比“隐藏”了负载耦合}
\label{sec:ds-hidecoupling}

\begin{insight}[$1/k_r^2$ 把构型相关惯量与关节耦合都压平]
\label{ins:ds-hide-coupling}
\cref{ch:arm-dyn} 反复强调：机械臂的惯量矩阵 $\mathbf{M}(\mathbf{q})$ \emph{强烈依赖构型}，对角元随位形变化可达数倍，非对角元（关节耦合）也不可忽略。但这是\emph{负载侧}的图景。从\emph{电机侧}看，由 \cref{eq:ds-Ieq-motor}，每个关节的有效惯量是 $I_m+M_{ii}(\mathbf{q})/k_r^2$：常数转子惯量 $I_m$ 占大头，构型相关的负载惯量被 $1/k_r^2$ 压成小扰动；同理，关节 $j$ 的运动经耦合项 $M_{ij}(\mathbf{q})$ 反映到关节 $i$ 的电机侧也被压 $1/k_r^2$。于是高减速臂的各关节\emph{近似解耦}，可当作一串\textbf{独立的、定常惯量的单关节二阶系统}来控（P4 \cref{sec:ctrl-robot} 的独立关节控制正基于此）。这就是“高减速比隐藏负载耦合”的精确含义。\textbf{反过来}：直驱臂（$k_r=1$）失去这层保护，构型相关惯量与耦合全部暴露，必须上计算力矩这类\emph{显式补偿动力学}的方法（\cref{sec:ctrl-robot-ctc}）。
\end{insight}

\begin{insight}[Paul 的实证：执行器惯量降低结构依赖]
\label{ins:ds-paul-evidence}
Paul 对一台二连杆臂逐数算过\cite{paul1981robot}：不计执行器时，关节 1 看见的有效惯量随 $\theta_2$ 与负载有 $3{:}1$ 乃至（外空间大负载时）$201{:}1$ 的剧烈变化；一旦把折算后的执行器惯量 $k_r^2 I_m$ 计入分母，“这些惯量变化的\emph{相对}幅度被显著压低、耦合惯量项的相对重要性也随之下降”。这与上一条洞见同义：执行器（经减速比放大的）惯量像一个“大基座项”，把构型相关部分稀释成扰动。它从\emph{数值}上印证了独立关节控制的可行性前提。
\end{insight}

\begin{pitfall}[“隐藏”不是“消失”——别把加速度盒当力矩约束]
减速比只是把耦合\emph{相对}压小，并未消除。在高速、大负载、低减速比或精密力控场景，被压下的耦合与构型相关惯量会重新抬头。一个常见且危险的近似是：以为“高减速 $\Rightarrow$ 惯量恒定 $\Rightarrow$ 固定加速度盒即代表力矩约束”。\cref{ch:arm-dyn} 已实证：有效惯量随构型变化，固定加速度盒在不同位形对应的力矩需求可差数倍（\cref{ch:arm-prac} 把加速度约束换成真力矩约束后轨迹快了数倍）。减速比给的是“可以\emph{先}用独立关节控制起步”的工程许可，不是“可以永远忽略动力学”的免责声明。
\end{pitfall}

% ════════════════════════════════════════════════════════════════
\section{谐波减速器、行星齿轮与关节柔性}
\label{sec:ds-flex}

上一节假定传动\emph{刚性无背隙}。真实减速器并非如此——尤其谐波减速器本就靠柔性变形工作。本节给出常用减速器的特性与“弹性关节”模型，为 \cref{sec:ds-bandwidth} 的共振分析铺路。

\subsection{两类主流减速器}
\label{sec:ds-reducers}

\paragraph{谐波减速器（harmonic drive / strain wave gear）。}\;由\emph{波发生器}（椭圆凸轮）、\emph{柔轮}（薄壁弹性外齿杯）、\emph{刚轮}（内齿圈）三件构成。波发生器把柔轮撑成椭圆，使其长轴两端的外齿与刚轮内齿啮合；波发生器每转一周，柔轮相对刚轮反向错过“柔轮与刚轮齿数之差”个齿。因齿差通常仅 $2$ 个，单级减速比可高达 $30$--$320$，且\emph{零背隙}（柔轮齿始终预压贴合）、同轴、轻量——这正是协作臂与精密关节的首选。代价是：柔轮靠\emph{弹性变形}传力，关节天生带一段\emph{扭转柔性}与轻微\emph{迟滞/运动学误差}（柔轮变形非完美正弦）。

\paragraph{行星齿轮（planetary gear）。}\;太阳轮（输入）带动数个行星轮，行星轮在固定齿圈内公转、带动行星架（输出）。多齿同时啮合使其\emph{承载大、刚度高、效率高}（$>0.9$），但\emph{有背隙}（齿侧间隙），需精密加工或预紧消隙；单级减速比一般 $3$--$10$，大减速比靠多级串联。重载工业臂常用。

\begin{note}[直驱与减速的取舍，一句话]
减速比越高，越能“隐藏”负载耦合、放大力矩、降低对控制的要求（\cref{ins:ds-hide-coupling}），但越容易引入柔性与（行星齿轮的）背隙，把结构共振频率\emph{拉低}（见 \cref{sec:ds-bandwidth}）。直驱反之：无背隙无传动柔性、共振频率高、带宽潜力大，但耦合全暴露、需复杂控制、电机大而贵。谐波减速器是个折中：高减速比$+$零背隙，但柔轮柔性是其固有“软肋”。
\end{note}

\subsection{弹性关节的二自由度模型}
\label{sec:ds-elasticjoint}

当传动柔性不可忽略时，电机角 $q_m$ 与连杆角 $q_\ell$ 不再差一个固定 $k_r$，而是被一段\emph{扭转弹簧} $k$（折算到连杆侧的等效刚度）隔开，成为\emph{两个独立自由度}。以单连杆为例，动能 $K=\tfrac12 J_\ell\dot q_\ell^2+\tfrac12 J_m\dot q_m^2$，势能含弹簧能 $P=Mg\ell(1-\cos q_\ell)+\tfrac12 k(q_\ell-q_m)^2$。代入欧拉--拉格朗日方程（暂略阻尼）得 Spong 的\textbf{弹性关节模型}\cite{spong1987flexible}：

\begin{equation}
\boxed{
\begin{aligned}
J_\ell\,\ddot q_\ell + Mg\ell\sin q_\ell + k\,(q_\ell-q_m) &= 0,\\
J_m\,\ddot q_m + k\,(q_m-q_\ell) &= u.
\end{aligned}}
\label{eq:ds-elastic}
\end{equation}

物理图像：电机力矩 $u$ \emph{先}推动电机惯量 $J_m$、经弹簧 $k$ \emph{再}传给连杆惯量 $J_\ell$。两个惯量被一根弹簧连着，构成一个二阶振子——这就是结构共振的来源。刚性极限 $k\to\infty$ 时 $q_m\to q_\ell$，二式相加退化回 \cref{eq:ds-spong-eom} 的刚性单连杆。

\begin{insight}[柔性把系统阶数抬高、把共振峰带进来]
\label{ins:ds-flex-order}
刚性关节是二阶（一对极点）；弹性关节是\emph{四阶}（\cref{eq:ds-elastic} 是两个耦合的二阶方程），多出一对靠近虚轴的\emph{欠阻尼共振极点}，频率约 $\sqrt{k/J_{\rm eff}}$。控制设计若沿用刚性二阶模型（即把这对极点当“未建模动态”忽略），就必须\emph{回避}激励它——这正是 \cref{sec:ds-bandwidth} 带宽上限的根源。若想把带宽推到共振频率之上，则\emph{必须}显式建模柔性（用 \cref{eq:ds-elastic} 这类四阶模型设计），控制随之复杂得多（奇异摄动、反步、基于无源性的方法等），已超出工业常规、属研究前沿。
\end{insight}

\begin{pitfall}[背隙：一段“死区$+$冲击”的非线性，比柔性更难缠]
柔性是\emph{线性}弹簧，尚可纳入线性模型；背隙（行星齿轮齿侧间隙、磨损松动）却是\emph{硬非线性}——电机反向时先空走一段“死区”（力矩为零、关节不动），齿重新咬合瞬间又是\emph{冲击}。后果：极限环振荡、过零畸变、定位重复性下降、测得的电机侧位置与真实关节位置在换向处错开一个间隙。谐波减速器零背隙的价值正在此。背隙无法靠线性控制根除，工程上靠机械预紧消隙、或在控制端做死区补偿/双编码器（电机侧$+$关节侧）闭环。
\end{pitfall}

% ════════════════════════════════════════════════════════════════
\section{液压执行器要点}
\label{sec:ds-hydraulic}

重载机械臂（搬运重物、大型工程臂）常用\textbf{液压伺服电机}：靠压缩流体推动活塞做容积变化而出力。线性液压缸（单活塞，行程有限）驱动移动关节；旋转液压马达（多活塞，行程无限）驱动转动关节。其静/动态性能可与电气伺服相当。

\subsection{输入/输出模型}
\label{sec:ds-hyd-model}

液压驱动的建模围绕三组关系：流量--压力平衡、流体与运动件的关系、运动件的机械平衡。设伺服阀供给的体积流量为 $Q$，则\textbf{流量平衡}

\begin{equation}
Q = Q_m + Q_l + Q_c,
\label{eq:ds-hyd-flow}
\end{equation}
$Q_m$ 传给马达、$Q_l=k_l P$ 是泄漏流量、$Q_c=\gamma V s P$ 是流体可压缩性引起的流量（高压 $\sim100\,$atm 下不可忽略，$P$ 为负载差压，$V$ 为瞬时流体体积，$\gamma$ 为压缩系数）。传给马达的流量正比于角速度 $Q_m=k_q\Omega_m$；运动件机械平衡同 \cref{eq:ds-mechanical}；驱动力矩正比于负载差压 $C_m=k_t P$。伺服阀阀芯位移对控制电压的传递函数 $\frac{X}{V_c}=\frac{G_s}{1+sT_s}$（$T_s\sim$ ms 可略），分配器线性化后 $P=k_x X-k_r Q$。

\begin{insight}[与电气驱动器形式类比，但压力反馈是“结构性”的]
\label{ins:ds-hyd-analogy}
把液压驱动的方块图与电气驱动（\cref{sec:ds-drives}）并排看，二者动态行为\emph{形式类似}：压力反馈环对应电流反馈环，压缩性时间常数对应电气时间常数。\textbf{但有一处本质不同}：液压里的压力反馈环是系统的\emph{结构性}特征（流体可压缩与泄漏天然存在），\emph{无法}像电气那样靠选 $C_i(s)$ 自由把驱动器配成速度型或力矩型——要改变它，必须额外加传感器与控制电路。所以“速度型/力矩型可选”是电气驱动器的特权。
\end{insight}

\paragraph{为何重载偏爱液压。}\;液压伺服\emph{功率重量比极高}、低速能直接出大力矩（常\emph{无需减速器}，从而无传动柔性/背隙）、静止时\emph{不烧毁}（电气电机靠重力静态堵转时电流持续发热、需抱闸）、自润滑且循环流体利于散热、在易燃环境本质安全。代价：需液压站、成本高难微型化、效率低、需维护、漏油污染。重载场合下液压站成本占比反而不高，故综合最优。气动则因流体可压缩误差难精确控制，多只用于夹爪开合。

% ════════════════════════════════════════════════════════════════
\section{传感：本体感知与外感知}
\label{sec:ds-sensors}

控制要闭环，先得\emph{测得到}。传感器分两类：\textbf{本体感知}（proprioceptive，测机械臂内部状态：关节位置、速度、力矩）与\textbf{外感知}（exteroceptive，测环境：力、距离、视觉）。本节概述与关节控制最相关者；视觉详见 \cref{ch:visual-servoing}，力控用的力传感器与 \cref{ch:arm-collision-safety} 的残差观测器互补。

\subsection{位置传感：编码器与旋变}
\label{sec:ds-encoders}

机器人几乎都用\emph{角位移}传感器（即便移动关节，伺服电机也是旋转型）。最常见的是\textbf{编码器}与\textbf{旋转变压器}。

\paragraph{绝对编码器。}\;光学玻璃盘上多条同心码道，每道按透/不透扇区编码；码道数 $=$ 字长 $=$ 分辨率。为避免多位同时跳变引发误读，用\textbf{格雷码}（相邻码字只变一位）。关节控制典型用 $\ge12$ 位（$1/4096$ 转分辨率）。绝对编码器断电不丢位置（位置直接编码在盘上）。有减速器时，关节侧一圈对应电机侧多圈，需电子电路计圈。

\paragraph{增量编码器。}\;两条正交（相位差 $90^\circ$）码道，既数转换次数又判转向；常加第三条“零位”道作机械零参考。比绝对编码器简单廉价、应用更广，但需计数电路把增量\emph{累加}成绝对位置，且位置存于易失存储，受扰动/掉电可能丢失。编码器自带信号处理电路，直接输出数字位置；外接电路还可由脉冲序列\emph{重构}速度（脉冲频率法宜高速、采样周期法宜低速）。

\paragraph{旋转变压器（resolver）。}\;紧凑坚固的机电式传感器，基于两电路互感，无机械限位连续传角。转子绕组馈正弦激励，定子两正交绕组感应出幅值随转角变化的电压，经旋变--数字转换器（RDC）按反馈原理解出数字角度，分辨率可达 16 位。坚固耐恶劣环境，常用于高温/振动场合。

\begin{pitfall}[编码器测的是电机轴，不是关节——柔性/背隙处二者背离]
\label{pit:ds-motor-side}
绝大多数工业关节把编码器装在\emph{电机轴}（高速侧）：分辨率被减速比 $k_r$ 放大（电机转 $k_r$ 圈关节才转一圈），且便于换相。代价是位置环闭在\emph{电机侧}：当关节有柔性（\cref{eq:ds-elastic}），$q_\ell\ne q_m/k_r$，二者差一段扭转变形；有背隙时换向处更差一个间隙。后果是电机侧“站得很稳”而\emph{末端却在抖或偏}——因为反馈量根本没看见连杆的真实位置。高精度/柔顺关节因此用\textbf{双编码器}（电机侧$+$关节侧）：电机侧编码器闭快速内环、关节侧编码器闭精度外环。这也是 \cref{ch:arm-ctrl} 关节控制必须正视的“反馈量在哪一侧”的问题。
\end{pitfall}

\subsection{速度传感：测速机}
\label{sec:ds-tacho}

速度虽可由位置差分重构，但常宁愿\emph{直接}测。\textbf{直流测速机}是个永磁小直流发电机，输出电压正比于转速；因有换向器存在残余纹波（频率随转速变，难滤）。\textbf{交流测速机}用两正交定子绕组$+$杯形转子，一相馈恒幅正弦、另一相感应出幅值正比于转速的正弦，纹波频率为激励两倍、可滤除；杯形转子惯量低、无滑动接触。二者性能相当。

\subsection{力/力矩传感：应变片、轴力矩、腕力}
\label{sec:ds-force}

力/力矩测量通常归结为测\emph{应变}（力使弹性元件微变形）。基本元件是\textbf{应变片}：金属丝受应变改变电阻，接入惠斯通电桥（Wheatstone bridge）把电阻变化转成电压

\begin{equation}
V_o=\Big(\frac{R_2}{R_1+R_2}-\frac{R_s}{R_3+R_s}\Big)V_i.
\label{eq:ds-bridge}
\end{equation}
为补偿温漂，在相邻桥臂贴一片不受力的应变片；为增灵敏度，用一拉一压两片置于相邻桥臂。

\paragraph{轴力矩传感器。}\;在电机与关节之间插一段低扭转刚度、高弯曲刚度的空心轴，贴应变片测扭转应变，经滑环馈电与取信。它测的是\emph{减速器输出端\emph{下游}传给关节的力矩}，故\emph{不}等于方块图里的电机驱动力矩 $C_m$（未含惯性/摩擦/上游传动的贡献）。这正是 \cref{ins:ds-vel-vs-torque} 所说“真力矩闭环”的硬件——比电流推算力矩更免疫参数漂移。

\paragraph{腕力传感器。}\;装在腕部（末端连杆与末端执行器之间），测末端坐标系下力$+$力矩共六分量。由若干弹性元件（如\emph{马耳他十字}布置的四根矩形梁，对侧各贴一对应变片成八桥）测出 $\ge8$ 路应变，再经\textbf{标定矩阵}线性映射到六维力旋量（力在前、力矩在后，记号同全书）。因结构不可能完全解耦，独立应变数多于六。该传感器是力控、装配、与环境接触任务的核心。

\begin{note}[力传感 vs 无传感残差估计：两条互补路线]
直接力/力矩传感器精度高但加成本、加腕部质量、降带宽。一条互补路线是\emph{无外置传感器}的力估计——De Luca--Mattone 广义动量观测器\cite{deluca2003actuator}：仅由关节力矩与本体动力学，以一阶滤波收敛出外力矩残差 $\mathbf r\to\boldsymbol\tau_{\rm ext}$，无需测关节加速度。\cref{ch:arm-collision-safety} 以 $\mathbf r$ 为输入做碰撞检测，正是这条路线的应用。二者按精度/成本/带宽权衡选用。
\end{note}

\subsection{惯性测量（IMU）一瞥}
\label{sec:ds-imu}

定基串联臂多不必装 IMU（关节编码器已给全构型）；但浮动基平台（移动操作、人形、四足，见 \cref{ch:arm-kin} 提及的浮动基与 P7/P9）必须测基座姿态/角速度/比力。IMU 由三轴陀螺$+$三轴加速度计构成，输出角速度与比力，经积分/滤波（与 P1 估计部的 ESKF 等对接）给出基座位姿。其偏置与噪声建模、与编码器/视觉的融合属估计部主题，此处仅指明接口：\emph{关节传感给“臂相对基座”，IMU 给“基座相对世界”，二者串联才得末端在世界系的状态}。

% ════════════════════════════════════════════════════════════════
\section{控制带宽与结构共振}
\label{sec:ds-bandwidth}

本章收束于一个对控制器设计有\emph{硬约束}意义的结论：电机$+$减速器$+$负载这条带柔性的链有一个\textbf{结构共振峰}，它给闭环带宽（从而给控制增益）划了一条不可逾越的天花板。这条天花板，正是 \cref{ch:arm-ctrl} 调增益时第一个要撞的墙。

\subsection{单关节模型：从二阶到“被忽略的三阶”}
\label{sec:ds-singlejoint}

在 \cref{ins:ds-hide-coupling} 的许可下，把单关节当作\emph{定常有效惯量}的二阶系统。略去电机电感（\cref{sec:ds-balances} 的降阶），力矩型驱动器近似为纯力矩源，单关节模型即

\begin{equation}
\tau = J_{\rm eff}\,\ddot\theta + b_{\rm eff}\,\dot\theta,\qquad J_{\rm eff}=I_{\max}+k_r^2 I_m,
\label{eq:ds-singlejoint}
\end{equation}
取 $I_{\max}$（构型相关惯量的\emph{最大}值）以保证全程不欠阻尼\cite{craig2005introduction}。这是 P4 \cref{sec:ctrl-robot} PD$+$前馈与 \cref{sec:ctrl-lqr} LQR 的被控对象。\textbf{但}这是把柔性当“未建模动态”\emph{丢掉}后的简化——真实链含 \cref{sec:ds-flex} 的柔性，是更高阶的、带共振峰的系统。控制律基于二阶模型设计没错，\emph{前提是不去激励那个被丢掉的共振}。

\subsection{结构共振频率的估算}
\label{sec:ds-resfreq}

任何柔性环节都可近似为“等效刚度 $k$ $+$ 等效质量/惯量 $m$”的弹簧--质量振子，其自然频率

\begin{equation}
\omega_{\rm res}=\sqrt{\frac{k}{m_{\rm eff}}}.
\label{eq:ds-omegares}
\end{equation}
对轴/梁这类分布柔性，可用集总模型估等效质量（如把质量 $m$ 的梁等效为端部 $0.23m$ 的质点、把分布惯量 $I$ 等效为端部 $0.33I$）。典型工业臂结构共振在 $5$--$25\,$Hz；去掉减速/传动柔性的直驱臂可高至 $70\,$Hz。

\begin{example}[一根轴的未建模共振]
\label{ex:ds-shaft-res}
设一根（无质量）扭转刚度 $400\,\mathrm{N\cdot m/rad}$ 的轴驱动 $1\,\mathrm{kg\cdot m^2}$ 的转动惯量。若建模时忽略了轴的柔性，则这个被忽略的共振频率为 $\omega_{\rm res}=\sqrt{400/1}=20\,\mathrm{rad/s}$。控制器若想把闭环自然频率推到 $20\,$rad/s 附近，就会激起这根轴的扭振。
\end{example}

\subsection{安全裕度 $\omega_n\le\tfrac12\omega_{\rm res}$ 的来由与应用}
\label{sec:ds-margin}

\begin{insight}[为什么是“共振的一半”]
\label{ins:ds-half-rule}
我们的二阶设计模型（\cref{eq:ds-singlejoint}）\emph{看不见}共振极点。若闭环带宽 $\omega_n$ 逼近 $\omega_{\rm res}$，控制器在共振频率附近仍有可观增益，就会\emph{激励}这对欠阻尼极点：轻则末端振铃、定位整定慢，重则正反馈进共振、闭环\emph{失稳}。为留足频率间隔，使控制器在 $\omega_{\rm res}$ 处的增益已充分滚降、不致激振，工程经验律是
\begin{equation}
\boxed{\;\omega_n\le\tfrac12\,\omega_{\rm res}.\;}
\label{eq:ds-half-rule}
\end{equation}
直觉：把共振峰留在控制器“管不到”的高频区，让二阶模型在自己有效的低频区说了算。这条经验律\emph{并非}严格定理（确切裕度依赖阻尼比、控制结构、采样率与对振铃的容忍度），而是一条稳妥的拇指规则——它把“增益越大越快越准”（\cref{ch:arm-ctrl} 的诉求）硬生生卡在一个上限，从此\emph{增益不能无限开大}。
\end{insight}

\begin{insight}[这条天花板把整本控制部串了起来]
\label{ins:ds-bandwidth-ceiling}
\cref{eq:ds-half-rule} 的影响远不止单关节调参：
\begin{itemize}
  \item \textbf{对增益选择}：PD 的 $K_p$ 直接定闭环刚度/带宽 $\omega_n=\sqrt{K_p/J_{\rm eff}}$，故 \cref{eq:ds-half-rule} 等价于 $K_p\le\tfrac14 J_{\rm eff}\,\omega_{\rm res}^2$——一个由\emph{机械结构}（共振）和\emph{执行器}（有效惯量）共同决定的增益上限。临界阻尼则取 $K_d=2\sqrt{K_p J_{\rm eff}}$。
  \item \textbf{对采样率}：数字控制还要求采样率至少为最高待抑制频率（含共振）的两倍（奈奎斯特），故 $\omega_{\rm res}$ 也设了采样率下限——这把 \cref{ch:arm-ctrl} 的实现频率与本章的机械共振绑在一起（Paul 的解算率--共振约束同此理）。
  \item \textbf{对本体设计}：想要高带宽（快、准、硬），就得\emph{从机械上抬高} $\omega_{\rm res}$——更高刚度、更轻末端、更短力臂、谐波减速器的零背隙、乃至直驱。这把“控制性能”反向变成了“机械与执行器设计指标”，呼应 \cref{ins:ds-power-backward} 的倒推思想。
\end{itemize}
\end{insight}

\begin{example}[把规则用到一次增益整定]
\label{ex:ds-gain-tuning}
设某关节有效惯量 $J_{\rm eff}=1$（已含折算电机惯量），最低未建模共振 $\omega_{\rm res}=8\,\mathrm{rad/s}$。按 \cref{eq:ds-half-rule} 取 $\omega_n=4\,\mathrm{rad/s}$；则比例增益 $K_p=J_{\rm eff}\,\omega_n^2=16$，微分增益取临界阻尼 $K_d=2\sqrt{K_p J_{\rm eff}}=8$。这就是“在不激励共振的前提下，把刚度开到尽可能大”的标准做法——增益不是越大越好，而是顶到共振墙为止。
\end{example}

\begin{pitfall}[别用单关节模型的成功去担保多关节/直驱的稳定]
\cref{eq:ds-singlejoint} 的二阶简化与 \cref{eq:ds-half-rule} 的拇指规则，建立在\emph{高减速比$+$定常有效惯量$+$忽略柔性}三条假设上。一旦进入直驱（$k_r=1$，耦合暴露、构型相关惯量剧变）、高速大负载（耦合抬头）、或要求带宽高过共振（柔性必须显式建模），这套简化与规则都失效，必须回到 \cref{ch:arm-dyn} 的全动力学$+$计算力矩、或四阶柔性模型$+$高级控制。换言之：本章给的是“工程起步线”，不是“理论封顶线”。
\end{pitfall}

% ════════════════════════════════════════════════════════════════
\section{符号表}
\label{sec:ds-notation}

\begin{center}
\begin{longtable}{ll}
\toprule
符号 & 含义 \\
\midrule
\endhead
$V_a,\,I_a$ & 电枢电压、电枢电流 \\
$R_a,\,L_a$ & 电枢电阻、电枢电感 \\
$V_g$ & 反电动势（back-EMF） \\
$k_e\,(k_v)$ & 反电动势/电压常数（$V\!\cdot\! s/\mathrm{rad}$） \\
$k_t$ & 转矩常数（$\mathrm{N\!\cdot\! m/A}$；SI 下数值 $=k_e$） \\
$C_m,\,C_l$ & 电机驱动力矩、负载反作用力矩 \\
$\Omega_m,\,\Theta_m$ & 电机轴角速度、角位移（频域） \\
$I_m\,(J_m),\,F_m$ & 电机转子惯量、黏性摩擦系数 \\
$I,\,F\,(J_\ell,\,B_\ell)$ & 负载惯量、负载黏性摩擦 \\
$G_v,\,T_v$ & 功率放大器增益、时间常数 \\
$k_i,\,C_i(s)$ & 电流反馈系数、电流调节器 \\
$k_r\,(\eta,\,r)$ & 减速比（电机转/关节转，$\gg1$） \\
$J_{\rm eff}$ & 关节侧看见的有效惯量 $I+k_r^2 I_m$ \\
$I_{\rm eq},\,F_{\rm eq}$ & 电机侧等效惯量 $I_m+I/k_r^2$、等效摩擦 \\
$k$ & 关节（传动）扭转刚度 \\
$q_m,\,q_\ell$ & 弹性关节的电机角、连杆角 \\
$M(s)$ & 控制电压$\to$电机转角传递函数 \\
$k_m,\,T_m$ & $M(s)$ 的等效增益、机电时间常数 \\
$Q,\,P$ & 液压流量、负载差压 \\
$\omega_{\rm res}$ & 最低结构共振频率 $\sqrt{k/m_{\rm eff}}$ \\
$\omega_n,\,\zeta$ & 闭环自然频率、阻尼比 \\
$K_p,\,K_d$ & 关节 PD 比例/微分增益 \\
\bottomrule
\end{longtable}
\end{center}

% ════════════════════════════════════════════════════════════════
\section{速查卡}
\label{sec:ds-cheatsheet}

\begin{center}
\begin{longtable}{p{0.30\linewidth}p{0.62\linewidth}}
\toprule
要点 & 一句话 \\
\midrule
\endhead
直流电机三式 & 电枢 $V_a=(R_a+sL_a)I_a+k_e\Omega_m$；转矩 $C_m=k_tI_a$；机械 $C_m=(sI_m+F_m)\Omega_m+C_l$。 \\
统一传递函数 & $M(s)=\dfrac{k_m}{s(1+sT_m)}$，电机是个积分器（转速积分得转角）。 \\
速度型驱动器 & $k_i=0$，$\omega_m\!\propto\!v_c'$，扰动系数 $R_a/k_ek_t$ 小、抑扰好 $\Rightarrow$ 独立关节控制用。 \\
力矩型驱动器 & $Kk_i\gg R_a$，力矩与转速解耦、$\propto$ 指令，抑扰差 $\Rightarrow$ 计算力矩/力控/RL 用。 \\
反映律 & 负载惯量/摩擦 $\div k_r^2$（电机侧），反作用力矩 $\div k_r$；平方源于“加速度换算$\times$力臂比”各一次。 \\
有效惯量 & 关节侧 $J_{\rm eff}=I+k_r^2 I_m$；高 $k_r$ 时转子惯量占大头、构型相关项被压平。 \\
隐藏耦合 & 高减速 $\Rightarrow$ 各关节近似解耦的定常二阶系统（独立关节控制的前提）；直驱则耦合全暴露。 \\
弹性关节 & 柔性把系统抬成四阶、带欠阻尼共振极点 $\sqrt{k/J_{\rm eff}}$；谐波减速器零背隙但柔轮有柔性。 \\
背隙 & 死区$+$冲击的硬非线性 $\Rightarrow$ 极限环、过零畸变；靠预紧/双编码器/死区补偿。 \\
液压 & 高功率重量比、低速大力矩、静止不烧毁；压力反馈是结构性的、不能自由配速度/力矩型。 \\
编码器在电机侧 & 测的是 $q_m$ 不是 $q_\ell$；柔性/背隙处二者背离 $\Rightarrow$ 高精度用双编码器。 \\
结构共振 & $\omega_{\rm res}=\sqrt{k/m_{\rm eff}}$，工业臂 $5$--$25\,$Hz，直驱可至 $70\,$Hz。 \\
带宽天花板 & $\omega_n\le\tfrac12\omega_{\rm res}\Rightarrow K_p\le\tfrac14 J_{\rm eff}\omega_{\rm res}^2$；增益顶到共振墙为止，不能无限开大。 \\
\bottomrule
\end{longtable}
\end{center}

% ════════════════════════════════════════════════════════════════
\section{练习}
\label{sec:ds-exercises}

\begin{practice}[本章练习]
\begin{enumerate}
  \item \textbf{[电机模型]} 由 \cref{eq:ds-electrical,eq:ds-mechanical,eq:ds-torque,eq:ds-backemf} 完整消元，推出“电枢电压 $V_a\to$ 电机角速度 $\Omega_m$”的二阶传递函数（不降阶，保留 $L_a$）。再令 $L_a\to0$、$F_m\ll k_ek_t/R_a$，验证退化为 \cref{eq:ds-Ms} 的一阶（对转速）形式。指出两个时间常数各对应什么物理过程。
  \item \textbf{[速度型 vs 力矩型]} 写出速度型 \cref{eq:ds-velgen-tf} 与力矩型 \cref{eq:ds-torquegen-tf} 对负载扰动 $C_l$ 的稳态增益，证明速度型的扰动增益 $R_a/k_ek_t$ 远小于力矩型的 $1/F_m$。据此解释：若上层用计算力矩控制（\cref{sec:ctrl-robot-ctc}），为何\emph{宁可}选抑扰差的力矩型？（提示：抑扰责任归谁。）
  \item \textbf{[反映律]} 某关节负载惯量随构型在 $I\in[2,6]\,\mathrm{kg\cdot m^2}$ 变化，转子惯量 $I_m=0.01\,\mathrm{kg\cdot m^2}$，减速比 $k_r=30$。求关节侧有效惯量 $J_{\rm eff}=I+k_r^2 I_m$ 的最小、最大值，并算出“构型引起的惯量变化”占有效惯量的百分比。把 $k_r$ 降到 $1$（直驱）重算，对比说明减速比如何“压平”惯量变化。
  \item \textbf{[弹性关节]} 对 \cref{eq:ds-elastic} 在平衡点 $q_\ell=q_m=0$ 线性化（小角度 $\sin q_\ell\approx q_\ell$，设 $Mg\ell=0$ 先看纯传动柔性），求系统的四个极点，证明其中一对是频率约 $\sqrt{k(J_m+J_\ell)/(J_mJ_\ell)}$ 的欠阻尼共振极点。讨论 $k\to\infty$ 时极点的去向。
  \item \textbf{[共振与带宽]} 一根扭转刚度 $800\,\mathrm{N\cdot m/rad}$、无质量的轴驱动 $1.2\,\mathrm{kg\cdot m^2}$ 的连杆（直驱）。求未建模共振 $\omega_{\rm res}$。按 \cref{eq:ds-half-rule} 给出允许的最大闭环 $\omega_n$，并算临界阻尼下的 $K_p,K_d$。若该轴改经 $k_r=10$ 的刚性齿轮驱动同一连杆，定性说明共振频率与允许带宽会怎样变（提示：有效惯量随 $k_r^2$ 增大）。
  \item \textbf{[传感]} 解释为何把编码器装在电机侧（高速侧）能提高有效角分辨率，但在有柔性的关节上会让“电机站稳而末端偏抖”。设计一个双编码器方案，说明电机侧与关节侧编码器各闭哪一个环、为什么这样分工（联系 \cref{pit:ds-motor-side} 与 \cref{ch:arm-ctrl} 的反馈量选择）。
  \item \textbf{[综合·选型倒推]} 给定一个关节任务：峰值关节力矩 $40\,\mathrm{N\cdot m}$、峰值关节角速度 $3\,\mathrm{rad/s}$、要求末端定位带宽 $\ge6\,$Hz。按 \cref{ins:ds-power-backward} 的倒推思路，依次估算：所需关节机械功率 $P_u$；若电机额定转速 $3000\,$rpm，需要的减速比量级；为达 $6\,$Hz 带宽对结构共振 $\omega_{\rm res}$ 的下限要求（用 \cref{eq:ds-half-rule}）；据此论证应选谐波减速器（零背隙、高刚度）还是普通行星齿轮。
\end{enumerate}
\end{practice}

% ════════════════════════════════════════════════════════════════
\section{延伸阅读}
\label{sec:ds-further}

\begin{itemize}
  \item \textbf{执行器与传感系统}：Siciliano, Sciavicco, Villani, Oriolo, \emph{Robotics: Modelling, Planning and Control}\cite{siciliano2009robotics} 第 5 章——本章驱动器方块图、速度/力矩型推导、液压模型、传感概览的主线源头，含旋变 RDC、腕力标定矩阵等工程细节。
  \item \textbf{单关节建模与结构共振带宽}：Craig, \emph{Introduction to Robotics}\cite{craig2005introduction} §9.9——有效惯量 $I+\eta^2 I_m$、$\omega_n\le\tfrac12\omega_{\rm res}$ 经验律、梁/轴集总模型估共振的出处。
  \item \textbf{弹性关节建模与控制}：Spong, \emph{Modeling and Control of Elastic Joint Robots}\cite{spong1987flexible} 与 \emph{Robot Modeling and Control}\cite{spong2006robot} §6.1——弹性关节二自由度模型 \cref{eq:ds-elastic}、奇异摄动与积分流形控制法。
  \item \textbf{有效惯量与耦合的数值实证}：Paul, \emph{Robot Manipulators}\cite{paul1981robot} §6.2--6.3——二连杆有效惯量随构型/负载剧变、执行器惯量降低结构依赖的逐数算例。
  \item \textbf{无传感力估计}：De Luca, Mattone, 广义动量观测器\cite{deluca2003actuator}——力矩型驱动器上做碰撞检测/外力估计的无外置传感器路线（\cref{ch:arm-collision-safety} 应用）。
\end{itemize}

% ════════════════════════════════════════════════════════════════
\section{后续关系}
\label{sec:ds-relations}

\begin{itemize}
  \item \textbf{承上 · 动力学}：本章把 \cref{ch:arm-dyn} 的负载方程 $\mathbf{M}(\mathbf{q})\ddot{\mathbf{q}}+\mathbf{C}\dot{\mathbf{q}}+\mathbf{g}=\boldsymbol{\tau}$ 右边的 $\boldsymbol{\tau}$ 落实为“电机经减速器产生”，并解释了高减速比为何让构型相关的 $\mathbf{M}(\mathbf{q})$ 在电机侧近似常数（\cref{ins:ds-hide-coupling}）。
  \item \textbf{启下 · 关节控制}：\cref{ch:arm-ctrl} 的独立关节控制以本章单关节模型 \cref{eq:ds-singlejoint} 为被控对象，其增益上限直接由本章带宽天花板 \cref{eq:ds-half-rule} 给出（\cref{ins:ds-bandwidth-ceiling}）；计算力矩控制（\cref{sec:ctrl-robot-ctc}）则要求本章的力矩型驱动器（\cref{ins:ds-vel-vs-torque}）。
  \item \textbf{对接 · 控制导论}：本章末端的二阶/四阶被控对象正是 P4 \cref{sec:ctrl-robot} PD$+$前馈与 \cref{sec:ctrl-lqr} LQR 的对象；那里讲“控制律怎么设计”，本章讲“对象是什么、增益顶到哪”。
  \item \textbf{对接 · 力控与安全}：本章的轴力矩/腕力传感器与 \cref{ch:arm-collision-safety} 的无传感残差观测器互补，共同支撑力控与碰撞检测。
  \item \textbf{对接 · 实践}：\cref{ch:arm-prac} 把力矩约束塞进轨迹时间参数化，其“真力矩 vs 加速度盒”的实证正回应本章 \cref{ins:ds-power-backward} 的功率倒推与减速比反映。
\end{itemize}
